% Part: methods
% Chapter: induction
% Section: strong-induction

\documentclass[../../../include/open-logic-section]{subfiles}

\begin{document}

\olfileid{mth}{ind}{str}

\olsection{प्रबल विगमन}

आधी चर्चिलेल्या विगमन तत्त्वात $P(0)$ सिद्ध करतो आणि $P(k)$ असेल तर $P(k+1)$ असते हेही सिद्ध करतो. दुसऱ्या भागात $P(k)$ सत्य आहे असे गृहीत धरून त्या गृहीतकावरून $P(k+1)$ सिद्ध करतो. धन निर्देशांकासाठी याऐवजी $P(k-1)$ गृहीत धरून $P(k)$ सिद्ध करता येते. महत्त्वाचे असे की कोणत्याही संख्येपासून तिच्या उत्तरवर्ती संख्येपर्यंतचे हे अनुमान करता आले पाहिजे: एखाद्या संख्येच्या पूर्ववर्ती संख्येसाठी विधान सत्य आहे असे गृहीत धरून तिच्यासाठी ते सिद्ध करता आले पाहिजे.

विगमन तत्त्वाचा आणखी एक प्रकार आहे. त्यात केवळ $k$ च्या पूर्ववर्ती $k-1$ साठी विधान खरे आहे असे न मानता,~$k$ पेक्षा लहान असलेल्या सर्व नैसर्गिक संख्यांसाठी ते खरे आहे असे मानतो आणि त्यावरून~$k$ साठी ते प्रस्थापित करतो. यानेही सर्व~$n \in \Nat$ साठी $P(n)$ मिळते. कारण $P(0)$ प्रस्थापित केल्यावर~$1$ पेक्षा लहान सर्व नैसर्गिक संख्यांसाठी $P$ खरे आहे असे प्रस्थापित झाले आहे. प्रत्येक $l<k$ साठी $P(l)$ असेल तर $P(k)$ असते, हे माहीत असल्यास विशेषतः $k=1$ साठीही ते माहीत असते. म्हणून $P(1)$ असा निष्कर्ष काढता येतो. यामुळे $P(0)$ आणि $P(1)$ सिद्ध झाले; म्हणजे सर्व $l<2$ साठी $P(l)$. तसेच प्रत्येक $l<2$ साठी $P(l)$ असेल तर $P(2)$ असते, हे सशर्त विधानही उपलब्ध असल्यामुळे~$P(2)$ निष्कर्ष म्हणून मिळते; आणि असेच पुढे.

खरे तर “प्रत्येक $k$ साठी, प्रत्येक $l<k$ साठी $P(l)$ असेल तर $P(k)$ असते” हे सामान्य सशर्त विधान प्रस्थापित करता आले, तर $P(0)$ वेगळे प्रस्थापित करण्याची गरज नसते; ते यावरूनच मिळते. “प्रत्येक $l<k$ साठी $P(l)$” सारखे सामान्य विधान कोणतेही $l<k$ नसतील तेव्हा सत्य असते, हे आठवा. हे रिक्त संख्यापनाचे प्रकरण आहे: कोणतेही $A$ नसतील, तर “सर्व $A$ हे $B$ आहेत” हे सत्य असते; कोणताही $x$ हा~$!A(x)$ पूर्ण करत नसेल, तर $\lforall[x][(!A(x) \lif !B(x))]$ सत्य असते. या प्रकरणात औपचारिक रूप “$\lforall[l][(l < k \lif P(l))]$” असे असेल; आणि कोणतेही $l < k$ नसतील, तर ते सत्य असते. $k=0$ असेल तेव्हा नेमके हेच घडते: कोणतीही नैसर्गिक संख्या $l<0$ नाही. म्हणून $P$ काहीही असले, तरी “प्रत्येक $l<0$ साठी $P(l)$” सत्य असते. त्यामुळे “प्रत्येक $l<k$ साठी $P(l)$ असेल तर $P(k)$ असते” याची सिद्धता आपोआप~$P(0)$ प्रस्थापित करते.

~$k$ साठी विधान प्रस्थापित करताना केवळ $k-1$ साठीच्या विधानावर अवलंबून राहता येत नसेल, पण एक किंवा अधिक $l<k$ साठी ते सत्य असल्याचे गृहीतक आवश्यक असेल, तर हा प्रकार उपयुक्त ठरतो.

\end{document}
